Тема работы~--- разработка системы мониторинга микросервисов с элементами машинного обучения.
Цель~--- построить воспроизводимый путь от разнородной телеметрии до модели, которая обнаруживает аномалии и может быть поставлена рядом с существующим комплектом мониторинга (Prometheus, Zabbix).

Задачи, соответствующие этапам исследования:

\begin{enumerate}
  \item реализовать ETL-конвейер, приводящий бенчмарки и промышленную телеметрию к единой схеме и признаковому пространству;
  \item сравнить несколько алгоритмов обнаружения аномалий на одних и тех же данных и выбрать модель (или ансамбль);
  \item обучить выбранную модель и сохранить воспроизводимые артефакты;
  \item оформить лёгкий сервис (контейнер) с конфигурируемым источником метрик и каналом алертов.
\end{enumerate}

Актуальность следует из того, что микросервисные системы порождают большой поток метрик, а пороги <<вручную>> плохо переносятся между сервисами и нагрузками.
Научная часть работы~--- корректная постановка эксперимента (разделение ролей датасетов, временной сплит, метрики качества) и сравнение детекторов.
Практическая часть~--- программный конвейер и последующий inference-сервис.

Структура диссертации. Глава~1 описывает источники данных и ETL (этап выполнен).
Глава~2~--- сравнение моделей (текущий этап).
Глава~3~--- архитектура программного продукта.
В заключении формулируются результаты и направления развития.

Черновик начат в сентябре 2026~г. Официальный шаблон кафедры будет подставлен позже; содержание глав при этом сохраняется.
